\begin{enumerate}
\def\labelenumi{\arabic{enumi}.}
\setcounter{enumi}{0}
\item
  \textbf{Introduction}
\end{enumerate}

\subsection{Mission Overview and Motivation}

Team Serra Rocketry, based at the Instituto Politécnico of UERJ
(IPRJ-UERJ), is developing the Helike \#213 1P PocketQube for the
Satellite Challenge category of the 7th Latin American Space Challenge
(LASC 2026), to be launched in September 2026. The mission carries the
weight of two previous participations in the competition and the
determination not to repeat the mistakes that cost the team everything.

\subsubsection{The LASC 2025 Lesson}

At LASC 2025, Serra Rocketry flew two missions: the SR Couto rocket
(\#100) and the SR Coutinho satellite (\#261). Both were built with
self-funded resources, unique components, and no redundancy. One day
before launch, the onboard electronics began to fail — a burned BMP and
an unstable LoRa link. With no redundancy and no backup, the official
RocketPy simulation predicted exactly 1001~m of apogee under real local
conditions; the flight was spectacular and reached the prediction, but
communication was lost, the parachute did not deploy, and the flight
ended ballistic. Both payloads were lost.

The lesson crystallized: \emph{optimization without redundancy is a
trap}. The Helike is the materialization of this lesson.

\subsection{The Problem with Conventional Parachutes}

The LASC 2025 experience corroborates what the literature already
points out: parachutes are a critical failure point in miniaturized
satellites:

\begin{itemize}
\item \textbf{Single point of failure:} if the sequencer fails, the
  parachute does not open — total loss.
\item \textbf{Complexity:} requires an actuator (servo/pyrotechnics) +
  firing logic + energy at the right moment.
\item \textbf{Volume:} parachutes and deployment mechanisms consume
  critical volume in miniaturized satellites.
\item \textbf{Induced rotation:} asymmetric parachutes can induce
  unstable rotation.
\item \textbf{Lines:} risk of structural entanglement, especially in
  aerodynamic PocketQubes.
\end{itemize}

\subsection{The SRAB Solution}

The SRAB (Bioinspired Autorotative Recovery System) was born directly
from the pain of losing everything: \emph{``what if recovery simply...
worked? No servo, no pyrotechnics, no software decision — just
physics.''}

\begin{itemize}
\item \textbf{Passive:} no actuators, no deployment electronics, no
  mechanism at all.
\item \textbf{Inherent:} recovery works by aerodynamic principle — the
  physics does the work.
\item \textbf{Predictable:} behavior modelable by rigid-body dynamics.
\item \textbf{Compact:} wings stowed during flight, passive deployment
  at ejection.
\end{itemize}

\subsection{Bioinspired Innovation}

The SRAB is based on a mechanism that nature has refined over millions
of years: the \textbf{autorotation of samara seeds} (\emph{Acer
rubrum}, maple seeds). These seeds evolved an asymmetric geometry where
the center of mass sits at one end and the aerodynamic lift at the
other, inducing stable conical rotation during fall. The Leading-Edge
Vortex (LEV) generated by the rotation doubles the effective lift,
resulting in very low descent velocities for the mass of the object.

Adapting this biological solution to aerospace engineering for
PocketQubes is a step toward simpler, more reliable, and more elegant
missions.

\subsection{Design Philosophy: The Three-Filter Review}

The project is reviewed continuously through a \textbf{three-filter
philosophy} that links regulatory compliance, failure analysis, and
schedule discipline directly to the code and hardware being built:

\subsubsection{Filter 1 — Compliance (SCSM LASC 2026)}

Every design decision is checked against the Satellite Challenge
Standards Manual (SCSM, Edition 7, Revision 1):

\begin{itemize}
\item \textbf{PQB 7.1.9/7.1.10 — Two kill switches on the Z axis:} the
  electrical schematic implements two independent kill switches
  (SW1/SW2, SPDT) that keep the satellite powered off during
  integration inside the deployer. \emph{Status: design decided,
  physical validation pending.}
\item \textbf{ECS 9.1.3 — $\geq$4~h autonomy:} the power budget
  (Section 3) predicts $>$40~h autonomy, a $>$10$\times$ margin over
  the requirement. A bench test with 20\% margin is the remaining
  evidence.
\item \textbf{ECS 9.1.5/9.1.6 — RBF pin:} a third switch (SW3, RBF)
  cuts all energy when inserted and is removed only after integration.
  The RBF circuit is in series with the main power rail.
\item \textbf{REC 10.1.1 — Independence from the rocket:} the
  satellite has no fixed electrical connection to the launch vehicle;
  it powers on at ejection and starts telemetry immediately.
\item \textbf{REC 10.2.1 — Non-parachute recovery notification:} the
  SRAB is a non-parachute recovery system and requires notification to
  LASC. \emph{Status: sent.}
\item \textbf{ECS 9.2.4 — Legal frequencies:} LoRa at 915~MHz
  (Americas band, Brazil). Verification of the specific authorization
  is in progress.
\end{itemize}

\subsubsection{Filter 2 — Murphy (What Can Go Wrong)}

Each Single Point of Failure (SPoF) is identified and mitigated in
\emph{both} hardware and firmware:

\begin{itemize}
\item \textbf{SRAB deployment failure} (highest risk): passive design
  with no hinges; two drop-test campaigns validated deployment; the
  wing count provides geometric redundancy. The firmware FSM defines
  safe states (\texttt{SAFE}, \texttt{PRE-FLIGHT}, \texttt{BOOST},
  \texttt{COAST}, \texttt{APOGEE}, \texttt{DESCENT},
  \texttt{POST-LANDING}) so that no single software path can prevent
  recovery.
\item \textbf{Energy failure / brownout:} robust BMS with overcharge,
  overdischarge, short-circuit and overcurrent protection; autonomy
  margin test as evidence.
\item \textbf{LoRa communication loss:} local storage (SD + LittleFS
  fallback) preserves all flight data. LoRa is support, not mission.
\item \textbf{Scheduler blocking:} every task in the cooperative
  scheduler is non-blocking with timeout; a hardware watchdog (TWDT,
  5~s) resets the system if any task stalls.
\item \textbf{Logging failure:} SD card with automatic LittleFS
  fallback on the internal flash.
\end{itemize}

\subsubsection{Filter 3 — Timeline (Deadlines)}

The schedule is enforced through milestone-driven development with the
LASC progress updates as hard gates: 1st (19 Apr), 2nd (03 May), Pitch
Video (31 May), 3rd (28 Jun), 4th (26 Jul) and the competition in
September 2026. Every deliverable is mapped to an owner and an
evidence artifact (video, log, bench test) so that the filter reviews
are data-driven rather than opinion-driven.

\subsection{Mission Objectives}

\subsubsection{Primary Objective — SRAB Validation}

Demonstrate that the SRAB is a viable recovery mechanism for 1P
PocketQubes, with terminal velocity within the LASC regulatory range,
predictable behavior, and no critical single points of failure.

\subsubsection{Secondary Objective — Flight Telemetry}

Collect flight data (pressure, temperature, acceleration, rotation,
GPS position) and transmit in real time via LoRa, demonstrating remote
sensing capability during all flight phases.

\subsubsection{Support Objective — Post-Landing Location}

Use LoRa RSSI triangulation (3 beacons) + GPS to locate the satellite
on the ground after landing, complementing SRAB recovery.

\subsection{Project Specifications}

\begin{center}
{\small
\setlength{\tabcolsep}{6pt}
\begin{tabularx}{\linewidth}{>{\raggedright}p{5.2cm}X}
\toprule
\textbf{Parameter} & \textbf{Specification} \\
\midrule
Form factor & PocketQube 1P (50$\times$50$\times$50 mm) \\
Recovery & SRAB — 4 autorotative wings (passive, non-parachute) \\
Terminal velocity (predicted) & 10.05 m/s (SCSM-compliant) \\
Microcontroller & ESP32-C3 Super Mini (single-core RISC-V) \\
Communication & LoRa RFM95W, 915 MHz \\
Navigation & GPS NEO-8M \\
Sensors & BME280 (pressure/temp/humidity) + ICM-20602 (6-axis IMU) \\
Storage & SD card + LittleFS fallback \\
Safety & 2 kill switches + RBF pin \\
Autonomy & $>$40 h (Li-Ion 6000 mAh) \\
\bottomrule
\end{tabularx}
}
\end{center}
